\documentclass[12pt]{article}
\usepackage{../../resource/shortex}
\usepackage{mathtools,amssymb}
\usepackage{enumitem}
\usepackage{datetime}
\usepackage{fancyhdr}
\usepackage{ifthen}
\usepackage{pifont}
\usepackage{mathrsfs}
\usepackage[normalem]{ulem}
%\RequirePackage[usenames,dvipsnames]{color}


%% showing solutions 
\newboolean{showsols}
\setboolean{showsols}{false}
\newcommand{\showsolutions}{\setboolean{showsols}{true}}
\newlength{\solspace}
\setlength{\solspace}{16em}
\newcommand{\nosolspace}{\setlength{\solspace}{0em}}

\newcommand{\solution}[2]{\ifthenelse{\boolean{showsols}}{{\par\textcolor{red}{Rubric criterion: #1}\newline\textcolor{blue}{#2}}}{\vspace{\solspace}}}
\newcommand{\npforprint}{\ifthenelse{\boolean{showsols}}{}{\newpage}}
\newcommand{\vsforprint}[1]{\ifthenelse{\boolean{showsols}}{}{\vspace{#1}}}

%% grading
\newcommand{\grades}[1]{~\newline\noindent\framebox[\textwidth][c]{E\quad \quad \quad \quad S \quad \quad \quad \quad R \quad \quad \quad \quad N}}

\newcommand{\questiongrades}{\noindent\emph{Each item is marked $\checkmark$, $\checkmark\!-$, or \ding{55}:}\begin{center}\emph{$\checkmark$ = capable  $~~$ $\checkmark\!-$ = mostly capable  $~~$ \ding{55} = not yet}\end{center} }

\newcommand{\triAssessments}{$\checkmark\quad\checkmark\!-\quad$\ding{55}\quad}
\newcommand{\biAssessments}{$\checkmark\quad$\ding{55}\quad}

%% line for writing name and BUID
\newcommand{\nameBUID}{\noindent\textbf{Name:}\underline{\phantom{XXXXXXXXXXXXXXXXXXX}}\hfill\textbf{BUID:}\underline{\phantom{XXXXXXXXXXXXX}}\newline}


%% formatting 
%\allowdisplaybreaks[3]
\renewcommand{\headrulewidth}{.4mm} % header line width
\renewcommand{\arraystretch}{1.25}


%% semester label, read from the "Semester:" line of important_dates.md
%% (path is relative to a document in materials/ModuleX/<folder>/)
\ExplSyntaxOn
\tl_new:N \g_teaching_semester_tl
\ior_new:N \g_teaching_dates_ior
\file_if_exist:nTF { ../../instructor_inputs/important_dates.md }
  {
    \ior_open:Nn \g_teaching_dates_ior { ../../instructor_inputs/important_dates.md }
    \ior_str_map_inline:Nn \g_teaching_dates_ior
      {
        \regex_extract_once:nnNT { \A \s* Semester: \s* (.*?) \s* \Z } {#1} \l_tmpa_seq
          {
            \tl_gset:Nx \g_teaching_semester_tl { \seq_item:Nn \l_tmpa_seq { 2 } }
            \ior_map_break:
          }
      }
    \ior_close:N \g_teaching_dates_ior
  }
  {
    \msg_new:nnn { teaching } { no-dates }
      { ../../instructor_inputs/important_dates.md~not~found;~semester~header~left~blank. }
    \msg_warning:nn { teaching } { no-dates }
  }
\cs_new:Npn \semester { \tl_use:N \g_teaching_semester_tl }
\ExplSyntaxOff

%% page style
\pagestyle{fancy}
\fancyhf{}
\rhead{\semester}
\lhead{CAS MA 213: Basic Statistics and Probability}
%\rfoot{\thepage}

%--------------------
\begin{document}

\begin{center}
\textbf{\Large Deliverable 2 \\ Slides, R Code, and Data}
\end{center}

Your group will submit draft \textbf{slides} summarizing your results, the \textbf{R code} used to generate the visualizations and analyses, and the \textbf{data file} if it is needed to run your code.
The slides should include the following sections:

\begin{enumerate}
    \item Title Slide (Project title, group members, date)
    \item Data Source and Description
    \item Categorical Data Analysis (EDA and Inference)
    \item Numerical Data Analysis (EDA and Inference)
    \item Conclusion and Future Work
    \item References
    \item Roles and Responsibilities of Each Group Member
\end{enumerate}

\vspace{6pt}
\noindent \textbf{Data Description Requirements:}
\begin{itemize}
    \item Cite the data source and state what one row represents.
    \item Provide a brief description of the categorical and numerical variables you use.
    \item Indicate whether there is any missing data, and how you handled it.
    \item Indicate whether there are any outliers or unusual values.
\end{itemize}

\vspace{6pt}
\noindent \textbf{Categorical Data Section -- Please include:}
\begin{itemize}
    \item The research question and the categorical variable(s) of interest.
    \item At least one summary table (frequency table or contingency table).
    \item At least one categorical visualization (e.g. bar plot or pie chart) with clear labels and title.
    \item At least one hypothesis test, including:
    \begin{itemize}
        \item The null and alternative hypotheses.
        \item A check of the conditions for inference.
        \item The test statistic and p-value.
    \end{itemize}
    \item A conclusion in the context of the data.
\end{itemize}

\vspace{6pt}
\noindent \textbf{Numerical Data Section -- Please include:}
\begin{itemize}
    \item The research question and the numerical variable(s) of interest.
    \item At least one summary statistics table (mean, median, sd, IQR, etc.).
    \item At least one numerical visualization (e.g. bar plot, histogram, boxplot, or scatterplot) with clear labels and title.
    \item At least one hypothesis test, including:
    \begin{itemize}
        \item The null and alternative hypotheses.
        \item A check of the conditions for inference.
        \item The test statistic and p-value.
    \end{itemize}
    \item A conclusion in the context of the data.
\end{itemize}

\vspace{6pt}
\noindent \textbf{R Code and Data Submission:}
\begin{itemize}
    \item Submit the \textbf{R code} used to perform your categorical and numerical analyses and generate visualizations as a \texttt{.R} or \texttt{.Rmd} file, along with the \textbf{data file} if it is needed to run your code.
    \item The code should be well-commented to explain the steps taken in the analysis.
    \item Ensure that the code can be run without errors, assuming the data file is in the same directory.
    \item You should only use functions and packages that have been covered in class, tutorials, or labs, or are part of the standard R distribution.
    \item Do not use functions that directly perform hypothesis tests (e.g., \texttt{t.test()}, \texttt{chisq.test()}, \texttt{prop.test()}). Compute test statistics and p-values step by step. If you use ANOVA, you may use \texttt{aov()}, \texttt{anova()}, or \texttt{lm()}.
\end{itemize}
\end{document}
